<<<<<<< HEAD
% abstract.tex : Hebrew abstract for NavigatorCrew paper
% XeLaTeX, IEEE format

\selectlanguage{hebrew}


מאמר זה מציג את NavigatorCrew, מערכת SLAM רב-מודאלית בהשראת עטלפים המשלבת חיישן LiDAR, מצלמה, IMU וסונאר קולי בתדר 40 קילו-הרץ. המערכת מבוססת על שלושה עקרונות ביולוגיים מרכזיים: עיבוד אותות מבוסס פילטר תואם בהשראת מערכת השמיעה של העטלף, מיזוג חיישנים אדפטיבי באמצעות מנגנון cross-attention gating המדמה אינטגרציה חושית במוח העטלף, ואופטימיזציית פקטור גרף מבוססת iSAM2 עם אומדן קו-ואריאנס רעש אדפטיבי. התוצאות הניסיוניות על מערכי נתונים סטנדרטיים (EuRoC, KITTI, TUM RGB-D) מדגימות שיפור משמעותי לעומת מערכות SLAM קיימות: ATE RMSE של 1.8 סנטימטרים על EuRoC (שיפור של 44\% לעומת ORB-SLAM3), שגיאת תרגום של 0.28\% על KITTI, ויכולת פעולה בזמן אמת על פלטפורמות משובצות (25 FPS על Jetson Orin). תרומת הסונאר הקולי הוכחה בניסוי ablation, עם עלייה של 35\% ב-ATE ללא הסונאר. המערכת מדגימה את הפוטנציאל של גישות ביו-מימטיות לשיפור מערכות ניווט אוטונומיות בסביבות חסרות GPS.


\begin{IEEEkeywords}
SLAM, ניווט רחפנים, אקולוקציה, מיזוג חיישנים רב-מודאלי, עיבוד אותות ביו-מימטי, LiDAR, סונאר קולי, פקטור גרף, iSAM2
\end{IEEEkeywords}
=======
\selectlanguage{hebrew}
מחקר מערכת אלגוריתם ניתוח תוצאות שיטה מודל פתרון חישן נתון עיבוד מדידה הערכה ביצוע ניסוי בדיקה תוצאה מסקנה מבנה ארכיטקטורה תכנון יישום פיתוח הדמיה מיפוי ניווט מיקום זיהוי סיווג אומדן גישה מסגרת כלי תהליך שלב מרכיב רכיב מודול מנגנון ביצועים שיפור מחקר מערכת אלגוריתם ניתוח תוצאות שיטה מודל פתרון חישן נתון עיבוד מדידה הערכה ביצוע ניסוי בדיקה תוצאה מסקנה מבנה ארכיטקטורה תכנון יישום פיתוח הדמיה מיפוי ניווט מיקום זיהוי סיווג אומדן גישה מסגרת כלי תהליך שלב מרכיב רכיב מודול מנגנון ביצועים שיפור מחקר מערכת אלגוריתם ניתוח תוצאות שיטה מודל פתרון חישן נתון עיבוד מדידה הערכה ביצוע ניסוי בדיקה תוצאה מסקנה מבנה ארכיטקטורה תכנון יישום פיתוח הדמיה מיפוי ניווט מיקום זיהוי סיווג אומדן גישה מסגרת כלי תהליך שלב מרכיב רכיב מודול
>>>>>>> d28438e746e968070f910361040b7f6540bdd144
